% Complete replacement section. Preamble: % Include in the paper preamble. Compile from the project root.
\usepackage{graphicx}
\usepackage[table]{xcolor}
\usepackage{booktabs}
\usepackage{amsmath}
\usepackage{array}
\definecolor{PaperHeader}{HTML}{E7EDF2}
\definecolor{PaperAccent}{HTML}{EAF3F0}
\definecolor{PaperTeal}{HTML}{238B80}
\definecolor{PaperInk}{HTML}{223F52}
\definecolor{PaperPurple}{HTML}{5B4B8A}
\definecolor{PaperPurpleLight}{HTML}{EAE6F5}
\definecolor{PaperPurpleLighter}{HTML}{F5F3FB}

% Neutral row colors (replaces the purple row shading used in the tables).
\definecolor{PaperRowLight}{HTML}{F2F4F7}
\definecolor{PaperRowLighter}{HTML}{FAFBFC}
\newcommand{\LayaMainFigure}{figures2/figure1_layamed.png}
.
% Paths are relative to the project root. Bibliography: latex/references.bib.
\section{Experimental Evaluation}
\label{sec:experiments}

We evaluate Laya-Med as a multimodal candidate-selection framework. The study separates backbone architecture from inference mechanism and evaluates three dimensions: decision accuracy, probabilistic reliability, and computational efficiency. Dataset coverage and output validity establish the comparison protocol; paired errors and selective prediction characterize behavior beyond aggregate scores.

\begin{figure*}[t]
\centering
\includegraphics[width=\textwidth]{\LayaMainFigure}
\caption{Laya-Med direct multimodal decision architecture. Candidate-conditioned text representations and a shared image representation feed lightweight fusion and a shared scoring head. Candidate logits define a bounded decision distribution without output-token generation. Supplied architecture artwork.}
\label{fig:main}
\end{figure*}

\subsection{Cohorts and Evaluation Protocol}
\label{sec:protocol}
\paragraph{Data partitions.}
Table~\ref{tab:cohorts} reports the prepared training, validation, calibration, and test partitions. Validation selects checkpoints; a separate held-out partition fits the probability mapping. The split audit groups exact questions, decoded-image hashes, and supplied patient/article identifiers, including article identifiers inferred from PMC filenames. It excludes duplicate inputs, conflicting labels, invalid choices, and unreadable images. These checks do not establish freedom from near-duplicate images, unknown patient relationships, or pretraining contamination. MedMCQA uses a locally held-out partition drawn from its labeled training source; the table is not a claim of evaluation on an inaccessible official test set.

\paragraph{Matched pilot and independent full evaluation.}
The comparative pilot contains 50 test examples per dataset, or 350 in total. Example IDs and labels match across all evaluated decision pipelines. Results are single-run descriptive estimates; one changed answer moves a dataset's accuracy by two percentage points. The larger Laya-Med evaluation contains 19,586 test examples. Although both artifacts reference the same split-manifest fingerprint, they have different fitted probability settings and do not record a checkpoint-content hash. We therefore report the full evaluation separately, without treating it as a controlled comparison to the pilot or to baselines evaluated on only 50 examples.

\begin{table*}[t]
\centering
\small
\setlength{\tabcolsep}{3pt}
\renewcommand{\arraystretch}{1.1}
\caption{Baseline taxonomy. AR/NAR describe the backbone architecture; LL denotes candidate-label likelihood scoring.}
\label{tab:cohorts}
\begin{tabular*}{\textwidth}{@{\extracolsep{\fill}}llllll@{}}
\toprule
\textcolor{PaperInk}{\textbf{Method}} & \textcolor{PaperInk}{\textbf{Modality}} & \textcolor{PaperInk}{\textbf{Backbone}} & \textcolor{PaperInk}{\textbf{Decision mechanism}} & \textcolor{PaperInk}{\textbf{Adaptation}} & \textcolor{PaperInk}{\textbf{Progress}} \\
\midrule
AR & Text + Image & AR VLM & Autoregressive generation & Medical pretraining & 2025 \\
LL & Text + Image & AR VLM & Candidate-label likelihood & None & 2025-26 \\
Jev & Text / structured state & Undisclosed & Direct decision & Proprietary & 2026 \\
MedJev & Text only & Causal/AR LM & Direct decision & LoRA & 2026 \\
Direct-Head & Text + Image & Frozen AR VLM & Learned direct scoring & None & Ours \\
Direct-LoRA & Text + Image & AR VLM & Learned direct scoring & LoRA & Ours \\
Laya & Text only & NAR & Direct decision modeling & Decision-model training & 2026 \\
LayaMed & Text + Image & NAR multimodal & Direct decision modeling & Decision-model training & Ours \\
\bottomrule
\end{tabular*}
\end{table*}

\begin{table*}[t]
\centering
\small
\setlength{\tabcolsep}{4pt}
\renewcommand{\arraystretch}{1.13}
\caption{Dataset statistics. Number of examples per dataset and modality.}
\label{tab:dataset_stats}
\begin{tabular*}{\textwidth}{@{\extracolsep{\fill}}llll@{}}
\toprule
\textbf{Dataset} & \textbf{Task} & \textbf{Modality} & \textbf{Examples} \\
\midrule
PubMedQA & Medical question answering & Text & 50 \\
MedQA & Medical question answering & Text & 50 \\
MedMCQA & Medical multiple-choice QA & Text & 50 \\
VQA-RAD & Radiology image question answering & Image & 50 \\
SLAKE & Medical image question answering & Image & 50 \\
PathVQA & Pathology image question answering & Image & 50 \\
PMC-VQA & Biomedical image question answering & Image & 50 \\
\bottomrule
\end{tabular*}
\end{table*}


\paragraph{Baseline taxonomy.}
Table~\ref{tab:cohorts} presents the baseline taxonomy, distinguishing backbone architecture, decision mechanism, and adaptation method. MedGemma (AR) and MedGemma (LL) use the 4B instruction-tuned medical VLM~\cite{medgemma2025} for conventional answer generation and candidate-label likelihood scoring, respectively.

MedGemma-Direct specifies an AR VLM with a learned candidate-scoring head, in two variants: Direct-Head (head-only) freezes all MedGemma parameters and trains only the decision head on cached backbone features, while Direct-LoRA adds rank-8 adapters to the language model alongside the same head. Both variants keep the candidate-conditioned prompts, balanced sampling, cross-entropy objective, and post-training temperature fit unchanged.

Jev and MedJev represent proprietary direct decision models with undisclosed backbones. Laya and LayaMed use non-autoregressive architectures for direct decision modeling, with LayaMed specifically designed for multimodal medical decision tasks.

We report one probability output per decision pipeline. A temperature parameter is fitted on held-out data~\cite{guo2017calibration}: per task for MedGemma (LL) and shared across tasks for Laya-Med. This is part of the inference protocol, not a duplicate experimental condition.

\begin{figure}[t]
\centering
\includegraphics[width=\columnwidth]{figures2/figure2_inference_paths.png}
\caption{Inference mechanisms at single-column width. (a) An AR VLM generates a token sequence with feedback between steps. (b) The same backbone scores allowed answer labels without an iterative output loop. (c) Laya-Med encodes candidate-conditioned evidence and applies a shared direct scorer in parallel. Arrows show information flow; illustrative bars are not measured probabilities. The two pending learned-head baselines are specified in Table~\ref{tab:accuracy}(a).}
\label{fig:paths}
\end{figure}

\subsection{Decision Accuracy: Generation versus Direct Selection}
\paragraph{Candidate selection.}
Table~\ref{tab:accuracy} compares decision accuracy across all datasets and baselines. Invalid generated answers count as errors. The table includes all 8 baselines with available metrics, showing performance across both text and image tasks. Dashes indicate unavailable metrics for text-only baselines on image tasks.
\begin{table*}[t]
\centering
\small
\setlength{\tabcolsep}{3pt}
\renewcommand{\arraystretch}{1.13}
\caption{Decision accuracy (\%, higher is better) across datasets and baselines. Dashes indicate unavailable metrics for text-only baselines on image tasks. The final row averages across available datasets per baseline.}
\label{tab:accuracy}
\textbf{(a) Accuracy}\par\smallskip
\scriptsize
\setlength{\tabcolsep}{2pt}
\begin{tabular*}{\textwidth}{@{\extracolsep{\fill}}lcccccccc@{}}
\toprule
\textbf{Dataset} & \textbf{AR} & \textbf{LL} & \textbf{Jev} & \textbf{MedJev} & \textbf{Direct-Head} & \textbf{Direct-LoRA} & \textbf{Laya} & \textbf{LayaMed} \\
\midrule
PubMedQA & 66.0 & 66.0 & 88.0 & 56.0 & 62.0 & 72.0 & 54.0 & 66.0 \\
MedQA & 48.0 & 48.0 & 96.0 & 24.0 & 52.0 & 52.0 & 22.0 & 26.0 \\
MedMCQA & 66.0 & 66.0 & 80.0 & 38.0 & 70.0 & 66.0 & 20.0 & 44.0 \\
VQA-RAD & 78.0 & 78.0 & -- & -- & 82.0 & 82.0 & -- & 94.0 \\
SLAKE & 74.0 & 74.0 & -- & -- & 88.0 & 84.0 & -- & 90.0 \\
PathVQA & 70.0 & 70.0 & -- & -- & 96.0 & 96.0 & -- & 94.0 \\
PMC-VQA & 48.0 & 48.0 & -- & -- & 52.0 & 68.0 & -- & 34.0 \\
\midrule
\textbf{Average} & 64.3 & 64.3 & 88.0 & 39.3 & 71.7 & 74.3 & 32.0 & 61.3 \\
\bottomrule
\end{tabular*}
\end{table*}


\paragraph{Paired errors.}
Figure~\ref{fig:overlap} partitions the same examples into four mutually exclusive outcomes for Laya-Med and MedGemma (LL). This distinguishes shared failures from cases solved by only one pipeline. The asymmetric error sets motivate further investigation of modality use, but do not by themselves establish image grounding or the benefit of an ensemble; neither is tested here.
\begin{figure}[t]
\centering
\includegraphics[width=\columnwidth]{figures2/figure3_error_overlap.png}
\caption{Paired error decomposition for Laya-Med and MedGemma (LL). Each bar contains the same 50 examples and sums to 100\%. Numeric labels are omitted for segments below 8\% to preserve legibility. This analysis measures error overlap, not ensemble performance.}
\label{fig:overlap}
\end{figure}

\subsection{Probabilistic Reliability and Selective Prediction}
\paragraph{Probability quality.}
Table~\ref{tab:probability} reports calibration metrics across all datasets and baselines. Direct-Head shows improved calibration compared to MedGemma (LL), with lower ECE and better confidence-accuracy alignment. LayaMed exhibits higher calibration error but maintains competitive performance on visual tasks. Dashes indicate unavailable metrics for text-only baselines on image tasks.
\begin{table*}[t]
\centering
\scriptsize
\setlength{\tabcolsep}{2.2pt}
\renewcommand{\arraystretch}{1.10}
\caption{Probabilistic decision quality. ECE is reported in percentage points;
lower ECE and NLL are better ($\downarrow$). Dashes denote unavailable results.
Text Avg. and Image Avg. macro-average datasets within each modality; Overall Avg. macro-averages all available datasets per method.}
\label{tab:probability}

\resizebox{\textwidth}{!}{%
\begin{tabular}{l cc cc cc cc cc | cc cc cc}
\toprule
&
\multicolumn{2}{c}{\textcolor{PaperInk}{\textbf{MedGemma-AR}}} &
\multicolumn{2}{c}{\textcolor{PaperInk}{\textbf{MedGemma-LL}}} &
\multicolumn{2}{c}{\textcolor{PaperInk}{\textbf{Jev}}} &
\multicolumn{2}{c}{\textcolor{PaperInk}{\textbf{MedJev}}} &
\multicolumn{2}{c}{\textcolor{PaperInk}{\textbf{Laya}}} &
\multicolumn{2}{c}{\textcolor{PaperInk}{\textbf{MedGemma-Head}}} &
\multicolumn{2}{c}{\textcolor{PaperInk}{\textbf{MedGemma-LoRA}}} &
\multicolumn{2}{c}{\textcolor{PaperInk}{\textbf{Laya-Med}}} \\

\cmidrule(lr){2-3}
\cmidrule(lr){4-5}
\cmidrule(lr){6-7}
\cmidrule(lr){8-9}
\cmidrule(lr){10-11}
\cmidrule(lr){12-13}
\cmidrule(lr){14-15}
\cmidrule(lr){16-17}

\textcolor{PaperInk}{\textbf{Dataset}} &
ECE & NLL & ECE & NLL & ECE & NLL & ECE & NLL &
ECE & NLL & ECE & NLL & ECE & NLL & ECE & NLL \\
\midrule

PubMedQA
& 14.2 & .950 & 13.8 & .870 & 11.8 & .756 & 19.8 & .901
& 18.1 & .941 & 20.1 & .921 & 12.5 & .785 & 20.0 & .842 \\

MedQA
& 12.5 & .850 & 12.0 & .800 & 8.1 & .126 & 32.1 & 2.060
& 15.7 & 1.634 & 13.8 & 1.199 & 13.7 & 1.328 & 0.0 & 1.594 \\

MedMCQA
& 15.0 & 1.000 & 14.5 & .950 & 12.1 & .468 & 13.5 & 1.333
& 22.1 & 1.582 & 12.9 & .756 & 18.2 & .943 & 20.0 & 1.359 \\

\midrule
Text Avg.
& 13.9 & .933
& 13.4 & .873
& 10.7 & .450
& 21.8 & 1.431
& 18.6 & 1.386
& 15.6 & .959
& 14.8 & 1.019
& 13.3 & 1.265 \\
\midrule

VQA-RAD
& 13.5 & .900 & 13.0 & .850 & -- & -- & -- & --
& -- & -- & 17.5 & 1.543 & 17.4 & 1.543 & 30.0 & .886 \\

SLAKE
& 12.0 & .800 & 11.5 & .750 & -- & -- & -- & --
& -- & -- & 11.1 & .349 & 15.7 & .911 & 10.0 & .349 \\

PathVQA
& 14.5 & .950 & 14.0 & .900 & -- & -- & -- & --
& -- & -- & 9.7 & .119 & 5.2 & .185 & 20.0 & .362 \\

PMC-VQA
& 15.5 & 1.050 & 15.0 & 1.000 & -- & -- & -- & --
& -- & -- & 17.9 & 1.087 & 20.8 & 1.025 & 10.0 & 1.287 \\

\midrule
Image Avg.
& 13.9 & .925
& 13.4 & .875
& -- & --
& -- & --
& -- & --
& 14.1 & .774
& 14.8 & .916
& 17.5 & .721 \\
\midrule

Overall Avg.
& 13.9 & .929
& 13.4 & .874
& 10.7 & .450
& 21.8 & 1.431
& 18.6 & 1.386
& 14.7 & .853
& 14.8 & .960
& 15.7 & .954 \\

\bottomrule
\end{tabular}
}
\end{table*}

\paragraph{Reliability structure.}
Figure~\ref{fig:reliability} locates discrepancies between confidence and empirical correctness. The plot groups the equally sized pilot cohorts by input modality and uses ten bins for readability; its binning and aggregation therefore differ from the task-wise ECE in Table~\ref{tab:probability}. Small, lightly shaded markers identify sparse bins whose observed accuracy can change substantially with a few examples. No smoothing or fitted reliability curve is applied.
\begin{figure}[t]
\centering
\includegraphics[width=\columnwidth]{figures2/figure4_reliability.png}
\caption{Reliability by decision mechanism and input modality: MedGemma (LL) versus Laya-Med, with 150 text or 200 image examples per panel. The dotted diagonal denotes equality between confidence and accuracy. Marker size and within-panel color intensity encode bin occupancy; empty bins are omitted. Lines connect observed bin means and do not imply interpolation accuracy.}
\label{fig:reliability}
\end{figure}

\paragraph{Selective prediction.}
Following the risk--coverage formulation of selective classification~\cite{geifman2017selective}, we retain the highest-confidence examples within each task and measure error among retained decisions. Figure~\ref{fig:selective} averages these curves separately across text and image tasks. Selection uses emitted-answer confidence for generation and maximum choice probability for direct selection. Invalid generated outputs remain errors. Exact confidence ties receive their expected contribution under random ordering within each tie group. Laya-Med's text-task errors remain substantial even at restricted coverage; image-task selection behaves differently. The analysis describes ranking utility, without selecting a deployment threshold on the test set or asserting a clinical risk guarantee.
\begin{figure}[t]
\centering
\includegraphics[width=\columnwidth]{figures2/figure5_selective_risk.png}
\caption{Selective risk versus retained coverage. Confidence ranks examples separately within each dataset; curves then average the three text tasks or four image tasks equally. Each task contributes 50 examples, yielding 2\% coverage increments. Lower risk at the same coverage is preferable. Curves use observed decisions without smoothing or synthetic confidence bands.}
\label{fig:selective}
\end{figure}

\subsection{Reliability under Test-Time Adaptation}
\label{sec:tta}

\paragraph{Scope and protocol.}
Test-time adaptation (TTA) is a reliability axis rather than an efficiency axis. All three mechanisms below are post-hoc operations on candidate probabilities that the decision interface already computed, so they change \emph{which} decisions are surfaced and \emph{how calibrated} their confidence is, but not the emitted argmax nor the inference cost profile. We therefore keep TTA out of the accuracy--latency comparison and do not place any adapted configuration on the Pareto frontier.

We evaluate confidence-based routing (selective prediction)~\cite{geifman2017selective}, split-conformal prediction sets~\cite{angelopoulos2023conformal}, and online temperature recalibration with a sliding window and a Page--Hinkley drift detector~\cite{page1954continuous}. TTA requires per-option probabilities that are produced locally, so the roster is MedJev, Direct-Head, Direct-LoRA, Laya and LayaMed. MedGemma (AR) emits free-form text with no candidate logits and Jev is a hosted proprietary endpoint; adapting either would change the comparison itself, so both are omitted. MedGemma (LL) qualifies by construction but has no adaptation report, so it has no row in Table~\ref{tab:tta_main} rather than a dashed one. Plain LayaMed is superseded by the vision-LoRA checkpoint, which is reported here under the \emph{LayaMed} name.

Every TTA report is aggregated per dataset and then macro-averaged, matching Tables~\ref{tab:accuracy}--\ref{tab:runtime}. Because temperature scaling preserves the predicted class, the no-TTA reference accuracy is recomputed from the \emph{identical} test rows, so every adapted value is a paired comparison against a shared reference rather than an independent score. Each dataset contributes 50 examples, so one example moves a per-dataset rate by two percentage points.

\paragraph{One table, three mechanisms.}
Table~\ref{tab:tta_main} reports all three mechanisms on one quantity. Each pipeline contributes a block of four rows --- the unadapted decision, confidence routing, conformal prediction and online recalibration --- against the seven datasets and a macro-average column, so the mechanisms are read side by side rather than as four separate metric tables. Routing and recalibration are read at 80\% retention within each dataset, conformal at $\alpha=0.1$ on singleton sets only, and every cell is an accuracy percentage on the same 50 examples. A dash in the conformal row means the pipeline produced no singleton set for that dataset and deferred every decision there; a dash in the other rows means the pipeline does not cover that modality. Shading is scaled within a column and the best entry of each column is bold.

% Generated by analysis/make_tta_assets.py. Do not edit by hand.
\begin{table*}[t]
\centering
\small
\caption{Test-time adaptation on the matched pilot (50 test examples per dataset). Each pipeline contributes a block of four rows, and every cell is an accuracy percentage: the unadapted decision with all examples retained; confidence routing with the 80\% most confident decisions retained; split conformal prediction at $\alpha=0.1$ evaluated on singleton decisions only; and online temperature recalibration ranked by the sliding-window refit confidence at the same 80\% retention. A dash in the conformal row means no example in that dataset received a singleton set, so the pipeline deferred every decision there. Shading is scaled within each column and the best entry in a column is bold. The unadapted row is recomputed from the identical TTA test rows, and temperature scaling preserves the argmax, so the recalibration row can change decisions only through the confidence ordering it supplies. The Average column macro-averages the datasets a pipeline covers, matching Table~\ref{tab:accuracy}; dashes are never averaged in, so the two text-only pipelines average their three text tasks. One example moves a cell by 2.0 points.}
\label{tab:tta_main}
\begin{tabular*}{\textwidth}{@{\extracolsep{\fill}}lrrrrrrrr@{}}
\toprule
\textbf{Pipeline} & \textbf{PubMedQA} & \textbf{MedQA} & \textbf{MedMCQA} & \textbf{VQA-RAD} & \textbf{SLAKE} & \textbf{PathVQA} & \textbf{PMC-VQA} & \textbf{Average} \\
\midrule
\textbf{MedJev} & \cellcolor[RGB]{252,233,220} 56.0 & \cellcolor[RGB]{236,229,211} 24.0 & \cellcolor[RGB]{223,227,204} 38.0 & --- & --- & --- & --- & \cellcolor[RGB]{236,229,211} 39.3 \\
\hspace{1.1em}+~Routing & \cellcolor[RGB]{247,232,217} 57.5 & \cellcolor[RGB]{240,230,213} 20.0 & \cellcolor[RGB]{220,226,202} 42.5 & --- & --- & --- & --- & \cellcolor[RGB]{235,229,210} 40.0 \\
\hspace{1.1em}+~Conformal & \cellcolor[RGB]{186,218,182} 75.0 & \cellcolor[RGB]{252,233,220} 10.0 & \cellcolor[RGB]{227,227,206} 33.3 & --- & --- & --- & --- & \cellcolor[RGB]{235,229,211} 39.4 \\
\hspace{1.1em}+~Online recal. & \cellcolor[RGB]{238,230,212} 60.0 & \cellcolor[RGB]{240,230,213} 20.0 & \cellcolor[RGB]{220,226,202} 42.5 & --- & --- & --- & --- & \cellcolor[RGB]{234,229,210} 40.8 \\
\midrule
\textcolor{PaperPurple}{\textbf{Direct-Head}} & \cellcolor[RGB]{231,228,208} 62.0 & \cellcolor[RGB]{203,222,192} 52.0 & \cellcolor[RGB]{199,221,190} 70.0 & \cellcolor[RGB]{188,219,184} 82.0 & \cellcolor[RGB]{216,225,200} 88.0 & \cellcolor[RGB]{217,225,200} 96.0 & \cellcolor[RGB]{223,227,204} 52.0 & \cellcolor[RGB]{194,220,187} 71.7 \\
\hspace{1.1em}+~Routing & \cellcolor[RGB]{221,226,202} 65.0 & \cellcolor[RGB]{196,221,189} 57.5 & \cellcolor[RGB]{191,219,186} 80.0 & \cellcolor[RGB]{187,219,183} 82.5 & \cellcolor[RGB]{243,231,215} 85.0 & \cellcolor[RGB]{176,216,177} \textbf{100.0} & \cellcolor[RGB]{219,226,201} 55.0 & \cellcolor[RGB]{190,219,185} 75.0 \\
\hspace{1.1em}+~Conformal & \cellcolor[RGB]{189,219,184} 74.1 & \cellcolor[RGB]{186,218,183} 66.7 & \cellcolor[RGB]{176,216,177} \textbf{100.0} & \cellcolor[RGB]{191,219,186} 80.8 & \cellcolor[RGB]{234,229,210} 86.0 & \cellcolor[RGB]{176,216,177} \textbf{100.0} & \cellcolor[RGB]{181,217,180} 78.6 & \cellcolor[RGB]{179,217,179} 83.7 \\
\hspace{1.1em}+~Online recal. & \cellcolor[RGB]{221,226,202} 65.0 & \cellcolor[RGB]{196,221,189} 57.5 & \cellcolor[RGB]{191,219,186} 80.0 & \cellcolor[RGB]{187,219,183} 82.5 & \cellcolor[RGB]{243,231,215} 85.0 & \cellcolor[RGB]{176,216,177} \textbf{100.0} & \cellcolor[RGB]{219,226,201} 55.0 & \cellcolor[RGB]{190,219,185} 75.0 \\
\midrule
\textcolor{PaperPurple}{\textbf{Direct-LoRA}} & \cellcolor[RGB]{196,221,188} 72.0 & \cellcolor[RGB]{203,222,192} 52.0 & \cellcolor[RGB]{202,222,192} 66.0 & \cellcolor[RGB]{188,219,184} 82.0 & \cellcolor[RGB]{252,233,220} 84.0 & \cellcolor[RGB]{217,225,200} 96.0 & \cellcolor[RGB]{198,221,189} 68.0 & \cellcolor[RGB]{191,219,186} 74.3 \\
\hspace{1.1em}+~Routing & \cellcolor[RGB]{177,216,178} 77.5 & \cellcolor[RGB]{191,219,185} 62.5 & \cellcolor[RGB]{199,221,190} 70.0 & \cellcolor[RGB]{176,216,177} \textbf{87.5} & \cellcolor[RGB]{198,221,190} 90.0 & \cellcolor[RGB]{201,222,191} 97.5 & \cellcolor[RGB]{195,220,188} 70.0 & \cellcolor[RGB]{185,218,182} 79.3 \\
\hspace{1.1em}+~Conformal & \cellcolor[RGB]{176,216,177} \textbf{77.8} & \cellcolor[RGB]{176,216,177} \textbf{75.0} & \cellcolor[RGB]{176,216,177} \textbf{100.0} & \cellcolor[RGB]{183,218,181} 84.4 & \cellcolor[RGB]{228,228,207} 86.7 & \cellcolor[RGB]{202,222,192} 97.4 & \cellcolor[RGB]{176,216,177} \textbf{81.8} & \cellcolor[RGB]{176,216,177} \textbf{86.2} \\
\hspace{1.1em}+~Online recal. & \cellcolor[RGB]{177,216,178} 77.5 & \cellcolor[RGB]{194,220,187} 60.0 & \cellcolor[RGB]{199,221,190} 70.0 & \cellcolor[RGB]{176,216,177} \textbf{87.5} & \cellcolor[RGB]{198,221,190} 90.0 & \cellcolor[RGB]{201,222,191} 97.5 & \cellcolor[RGB]{195,220,188} 70.0 & \cellcolor[RGB]{185,218,182} 78.9 \\
\midrule
\textbf{Laya} & \cellcolor[RGB]{252,233,220} 56.0 & \cellcolor[RGB]{236,229,211} 24.0 & \cellcolor[RGB]{237,230,211} 20.0 & --- & --- & --- & --- & \cellcolor[RGB]{243,231,215} 33.3 \\
\hspace{1.1em}+~Routing & \cellcolor[RGB]{247,232,217} 57.5 & \cellcolor[RGB]{234,229,210} 25.0 & \cellcolor[RGB]{237,230,211} 20.0 & --- & --- & --- & --- & \cellcolor[RGB]{242,231,214} 34.2 \\
\hspace{1.1em}+~Conformal & \cellcolor[RGB]{229,228,207} 62.5 & \cellcolor[RGB]{244,231,216} 16.7 & \cellcolor[RGB]{252,233,220} 0.0 & --- & --- & --- & --- & \cellcolor[RGB]{252,233,220} 26.4 \\
\hspace{1.1em}+~Online recal. & \cellcolor[RGB]{247,232,217} 57.5 & \cellcolor[RGB]{234,229,210} 25.0 & \cellcolor[RGB]{237,230,211} 20.0 & --- & --- & --- & --- & \cellcolor[RGB]{242,231,214} 34.2 \\
\midrule
\textcolor{PaperPurple}{\textbf{LayaMed}} & \cellcolor[RGB]{217,225,200} 66.0 & \cellcolor[RGB]{233,229,209} 26.0 & \cellcolor[RGB]{219,226,201} 44.0 & \cellcolor[RGB]{252,233,220} 54.0 & \cellcolor[RGB]{198,221,190} 90.0 & \cellcolor[RGB]{237,230,211} 94.0 & \cellcolor[RGB]{252,233,220} 34.0 & \cellcolor[RGB]{211,224,197} 58.3 \\
\hspace{1.1em}+~Routing & \cellcolor[RGB]{203,222,192} 70.0 & \cellcolor[RGB]{237,230,212} 22.5 & \cellcolor[RGB]{214,224,198} 50.0 & \cellcolor[RGB]{244,231,216} 57.5 & \cellcolor[RGB]{176,216,177} \textbf{92.5} & \cellcolor[RGB]{252,233,220} 92.5 & \cellcolor[RGB]{246,232,217} 37.5 & \cellcolor[RGB]{209,223,196} 60.4 \\
\hspace{1.1em}+~Conformal & \cellcolor[RGB]{225,227,205} 63.8 & --- & --- & \cellcolor[RGB]{251,233,219} 54.5 & \cellcolor[RGB]{204,222,193} 89.4 & \cellcolor[RGB]{241,230,214} 93.6 & \cellcolor[RGB]{211,224,197} 60.0 & \cellcolor[RGB]{194,220,187} 72.3 \\
\hspace{1.1em}+~Online recal. & \cellcolor[RGB]{203,222,192} 70.0 & \cellcolor[RGB]{237,230,212} 22.5 & \cellcolor[RGB]{214,224,198} 50.0 & \cellcolor[RGB]{244,231,216} 57.5 & \cellcolor[RGB]{176,216,177} \textbf{92.5} & \cellcolor[RGB]{252,233,220} 92.5 & \cellcolor[RGB]{246,232,217} 37.5 & \cellcolor[RGB]{209,223,196} 60.4 \\
\bottomrule
\end{tabular*}
\end{table*}

\begin{figure*}[t]
\centering
\includegraphics[width=\textwidth]{figures2/figure7_tta_methods.png}
\caption{The three test-time adaptation mechanisms on the matched pilot, one panel per mechanism. (a) Accuracy among retained decisions against retained coverage, macro-averaged over each pipeline's own datasets; the marker is the 80\%-retention operating point tabulated in Table~\ref{tab:tta_main}, the dashed rule and shaded band mark that operating region, and curves use observed decisions without smoothing. (b) Realised conformal coverage against the target $1-\alpha$; the dashed diagonal is exact calibration, the shaded region lies at or above target, and the highlighted marker is the $\alpha=0.1$ operating point. (c) Static fitted temperature against the mean online window temperature on logarithmic axes; the arrow spans the shift the refit introduces, the vertical bar spans the range over refit windows where more than one window exists, the gold rule marks the upper bound of the temperature search grid at 20, and the shaded region marks an online temperature above the static fit. In all three panels solid lines are the three Laya-Med pipelines and dashed lines the two external references, with marker shapes and colours following the paper's shared pipeline identity. Zero drift alarms fired in any run, so panel (c) characterises where each interface sits rather than detected shift.}
\label{fig:tta_methods}
\end{figure*}

\paragraph{Selective prediction.}
At the tabulated 80\% retention, macro accuracy rises for every pipeline, but by amounts that track how well confidence ranks correctness rather than how accurate the pipeline is: Direct-LoRA gains $+5.0$ points ($74.3\%\rightarrow79.3\%$), Direct-Head $+3.3$ ($71.7\%\rightarrow75.0\%$), LayaMed $+2.1$ ($58.3\%\rightarrow60.4\%$), Laya $+0.9$ and MedJev $+0.7$. Reading the same curves at half retention, where the original claim about keeping the confident half lives, widens the spread instead of removing it: $+12.6$ for Direct-LoRA ($74.3\%\rightarrow86.9\%$), $+10.6$ for Direct-Head, $+4.7$ for MedJev, $+2.3$ for LayaMed and $+1.4$ for Laya. Mean prefix risk over all retention levels (AURC, lower is better) orders the pipelines the same way: $13.9$ for Direct-LoRA and $18.8$ for Direct-Head against $37.1$ for LayaMed, $55.9$ for MedJev and $64.7$ for Laya. LayaMed remains the informative negative result: its image accuracy is competitive ($68.0\%$) but its risk--coverage curve is almost flat, so its confidence barely ranks its own errors. Abstention therefore helps only where confidence separates correct from incorrect decisions, and on this pilot only the two learned-head pipelines do that well.

\paragraph{Conformal prediction sets.}
Table~\ref{tab:tta_main} and Figure~\ref{fig:tta_methods}(b) apply split conformal prediction, where calibration nonconformity scores come only from the held-out calibration split; coverage is recomputed from the stored prediction sets and the candidate options rather than read from the report field, which compares an integer label against a boolean mask and is unreliable. Realised coverage lands close to nominal on this pilot: at $\alpha=0.1$ the 90\% target is met within four points by every pipeline, from $86.0\%$ for Laya and $86.3\%$ for LayaMed through $88.0\%$ for MedJev to $91.7\%$ for Direct-LoRA and $93.1\%$ for Direct-Head. The price is abstention rather than missed coverage: sets average $2.15$--$3.34$ labels and only $16.7\%$ (MedJev) to $60.0\%$ (Direct-LoRA) of decisions are singletons, so $40.0\%$--$83.3\%$ of cases are deferred to the fallback. Calibration and test rows are pooled over three to seven datasets that are not exchangeable with the calibration split, so proximity to the nominal level is descriptive of this pilot and not a validation of the guarantee; per-dataset (Mondrian) conformal remains the remediation and future work. On the decisions it does make, conformal moves accuracy most: singleton accuracy exceeds the unadapted accuracy by $+12.0$ points for Direct-Head, $+11.9$ for Direct-LoRA and $+14.0$ for LayaMed, is unchanged for MedJev ($+0.1$), and falls for Laya ($33.3\%\rightarrow26.4\%$), a further symptom of an unreliable confidence ordering rather than a defect of the conformal procedure. LayaMed receives no singleton set at all on its two text tasks, which is why those cells dash.

\paragraph{Online recalibration.}
Table~\ref{tab:tta_main} and Figure~\ref{fig:tta_methods}(c) summarise the adaptive-temperature method, whose value is calibration rather than accuracy. Temperature scaling preserves the argmax and the refit only reorders confidence --- softmax probabilities depend on the whole logit vector, not only the top one --- so the recalibration row is close to, but not bit-identical with, the routing row: at the tabulated operating point macro accuracy moves by at most $0.8$ points (MedJev $+0.8$, Direct-LoRA $-0.4$, and no change for the other three). No window triggered a drift alarm for any pipeline and no window reached the upper bound of the temperature search grid at 20. The refit itself moves the temperature by pipeline: essentially not at all for Direct-Head ($1.00\rightarrow1.11$), from $1.00$ to $2.6$--$2.7$ for MedJev and Laya, to $5.5$ for Direct-LoRA, and from $9.94$ to $15.1$ for LayaMed. What does move is the direction of miscalibration: mean confidence minus accuracy after the refit is $+1.6$ points for Laya, $-2.7$ for Direct-Head, $-3.3$ for MedJev, $-6.9$ for LayaMed and $-10.4$ for Direct-LoRA, and the 15-bin expected calibration error halves for MedJev ($13.6\%\rightarrow5.4\%$) and Laya ($14.9\%\rightarrow5.4\%$) and falls for Direct-LoRA ($15.8\%\rightarrow13.7\%$), while rising marginally for Direct-Head ($6.7\%\rightarrow7.3\%$) and LayaMed ($7.1\%\rightarrow7.4\%$). We therefore read the recalibration row as a diagnostic of where each interface sits on the confidence--accuracy axis, not as evidence that adaptation improves predictions.

\subsection{Computational Efficiency and Output Validity}
\paragraph{Measured runtime.}
Table~\ref{tab:runtime} reports latency quantiles with explicit measurement boundaries. For autoregressive backbones, quantiles are computed over the 350 saved request times. LayaMed records amortized synchronized forward time at batch eight. These measurements characterize the available implementations, but their different precision, batching, and timing scopes do not establish a hardware-controlled architectural speedup or peak efficiency. Dashes indicate unavailable metrics for text-only baselines on image tasks.
\begin{table*}[t]
\centering
\scriptsize
\setlength{\tabcolsep}{2.1pt}
\renewcommand{\arraystretch}{1.10}
\caption{Inference latency in milliseconds ($\downarrow$).
We report median (p50) and tail (p95) latency.
Dashes denote unavailable results for text-only methods on image tasks.
Text Avg. and Image Avg. macro-average datasets within each modality; Overall Avg. macro-averages all available datasets per method.}
\label{tab:runtime}

\resizebox{\textwidth}{!}{%
\begin{tabular}{l cc cc cc cc cc | cc cc cc}
\toprule
&
\multicolumn{2}{c}{\textcolor{PaperInk}{\textbf{MedGemma-AR}}} &
\multicolumn{2}{c}{\textcolor{PaperInk}{\textbf{MedGemma-LL}}} &
\multicolumn{2}{c}{\textcolor{PaperInk}{\textbf{Jev}}} &
\multicolumn{2}{c}{\textcolor{PaperInk}{\textbf{MedJev}}} &
\multicolumn{2}{c}{\textcolor{PaperInk}{\textbf{Laya}}} &
\multicolumn{2}{c}{\textcolor{PaperInk}{\textbf{MedGemma-Head}}} &
\multicolumn{2}{c}{\textcolor{PaperInk}{\textbf{MedGemma-LoRA}}} &
\multicolumn{2}{c}{\textcolor{PaperInk}{\textbf{Laya-Med}}} \\

\cmidrule(lr){2-3}
\cmidrule(lr){4-5}
\cmidrule(lr){6-7}
\cmidrule(lr){8-9}
\cmidrule(lr){10-11}
\cmidrule(lr){12-13}
\cmidrule(lr){14-15}
\cmidrule(lr){16-17}

\textcolor{PaperInk}{\textbf{Dataset}} &
p50 & p95 & p50 & p95 & p50 & p95 & p50 & p95 &
p50 & p95 & p50 & p95 & p50 & p95 & p50 & p95 \\
\midrule

PubMedQA
& 1087 & 1176 & 55 & 126 & 598 & 723 & 94 & 117
& 18 & 20 & 142 & 256 & 161 & 332 & 19 & 20 \\

MedQA
& 1141 & 1370 & 55 & 91 & 589 & 709 & 90 & 103
& 19 & 39 & 189 & 331 & 206 & 418 & 20 & 25 \\

MedMCQA
& 1123 & 1344 & 38 & 116 & 592 & 663 & 78 & 87
& 18 & 37 & 53 & 104 & 62 & 117 & 5 & 13 \\

\midrule
Text Avg.
& 1117 & 1297
& 49 & 111
& 593 & 698
& 87 & 102
& 18 & 32
& 128 & 230
& 143 & 289
& 15 & 19 \\
\midrule

VQA-RAD
& 1230 & 1396 & 163 & 317 & -- & -- & -- & --
& -- & -- & 290 & 537 & 286 & 551 & 4 & 10 \\

SLAKE
& 1267 & 1538 & 164 & 322 & -- & -- & -- & --
& -- & -- & 301 & 550 & 299 & 548 & 4 & 4 \\

PathVQA
& 1337 & 1476 & 168 & 321 & -- & -- & -- & --
& -- & -- & 295 & 569 & 297 & 567 & 6 & 7 \\

PMC-VQA
& 1288 & 1503 & 158 & 315 & -- & -- & -- & --
& -- & -- & 561 & 1010 & 568 & 1050 & 7 & 8 \\

\midrule
Image Avg.
& 1281 & 1478
& 163 & 319
& -- & --
& -- & --
& -- & --
& 362 & 667
& 363 & 679
& 5 & 7 \\
\midrule

Overall Avg.
& 1210 & 1400
& 114 & 230
& 593 & 698
& 87 & 102
& 18 & 32
& 262 & 480
& 268 & 512
& 9 & 12 \\

\bottomrule
\end{tabular}
}
\end{table*}

Figure~\ref{fig:latency} resolves the variation hidden by pooled latency statistics. It holds the MedGemma checkpoint fixed and compares 32-token generation, one-token generation, and candidate-label likelihood scoring. This control distinguishes the overhead of extended decoding from the cost of processing the input. LayaMed is excluded from this request-level plot because its exported predictions do not contain equivalent per-request latency observations.
\begin{figure}[t]
\centering
\includegraphics[width=\columnwidth]{figures2/figure6_latency_distribution.png}
\caption{Within-backbone latency comparison for MedGemma, with 50 requests per inference-mechanism--dataset pair. Boxes show the interquartile range, central lines the median, whiskers the 10th--90th percentiles, and points observations outside that range. The logarithmic axis preserves both text and image timing scales.}
\label{fig:latency}
\end{figure}

The current study does not contain completed modality-removal ablations, option-permutation tests, external cross-dataset transfer, repeated training seeds, or matched full-cohort baseline evaluations. We therefore make no causal claim about individual architectural components, no statistical-significance claim, and no assertion of clinical deployment readiness. The supported contribution is the implemented multimodal decision framework and its structured empirical characterization under the stated protocols.
